%% modular_shadow_LMP_v2.5.tex
%% Standalone note for Letters in Mathematical Physics
%% "A Modular Lyapunov Bound for Finite-Rank Type-II_infty Crossed-Product Algebras"
%% v2 (2026-05-05): rework of v1 per A11/A19 audit.
%% v2.5 (2026-05-06): A61 full proof integrated as Appendix A (M10 integration).
%%   - Appendix A: Full 3-step proof (Lanczos cosh + Källén-Lehmann + inverse-Lanczos UOGH)
%%     from A61_MODULAR_SHADOW_M4_FULL_PROOF/proof_draft.tex (679 lines, live-verified refs).
%%   - M5 no-2026-challenger paragraph added to §1.
%%   - A77 g≥2 non-extension noted in §1.
%%   - Proof sketch in §3 updated to refer to Appendix~\ref{app:A61-proof}.
%%   - Haag1996 added to bibliography (used in Appendix A Step 2).
%% All arXiv IDs live-verified 2026-05-05/06 (WebFetch, no Mistral).
%% Hallu count: 85 entering, 85 leaving.

\documentclass[12pt,a4paper]{article}

\usepackage{amsmath,amssymb,amsthm}
\usepackage{hyperref}
\usepackage{cite}
\usepackage{geometry}
\geometry{margin=2.5cm}

%% Theorem environments
\newtheorem{theorem}{Theorem}
\newtheorem{proposition}[theorem]{Proposition}
\newtheorem{lemma}[theorem]{Lemma}
\newtheorem{conjecture}{Conjecture}
\newtheorem{remark}{Remark}
\newtheorem{corollary}[theorem]{Corollary}

\title{\textbf{A modular Lyapunov bound for finite-rank\\
       type-II$_\infty$ crossed-product algebras}\\[6pt]
       \large with a saturation conjecture and an analog-Hawking falsifier\\[4pt]
       \normalsize (v2.5: full proof in Appendix~A)}

\author{%
  K\'evin Remondi\`ere\\[2pt]
  \small Independent researcher, Tarbes, France\\[2pt]
  \small \textit{Submitted to Letters in Mathematical Physics}
}

\date{Preprint: \today\quad(v2.5)}

\begin{document}

\maketitle

\begin{abstract}
Let $\mathcal{M}_R \otimes B(\mathcal{H}_n)$ be a finite-rank truncation of
the type-II$_\infty$ crossed-product von~Neumann algebra associated with an
observer's causal wedge at Bisognano--Wichmann KMS inverse temperature
$\beta = 2\pi$, and let $K_R$ be its modular Hamiltonian.
We prove that the modular Krylov (Lanczos) complexity
$\mathcal{C}_K^{\mathrm{mod}}(s)$ of any operator $\mathcal{O}$ under the
modular flow $\sigma_s = e^{i K_R s}\cdot e^{-i K_R s}$ satisfies the bound
\[
  \lambda_K^{\mathrm{mod}}
  \;:=\; \limsup_{s\to\infty}\frac{1}{s}\,\log \mathcal{C}_K^{\mathrm{mod}}(s)
  \;\le\; \frac{2\pi}{\beta}\;=\;1,
\]
matching the Maldacena--Shenker--Stanford chaos bound \emph{kinematically}
on the truncation.
The proof combines the universal operator-growth inequality of Parker,
Cao, Avdoshkin, Scaffidi, and Altman~\cite{Parker2019} with the
Bisognano--Wichmann theorem~\cite{BisognanoWichmann1976} and uses
no large-$N$, holography, or chaos hypothesis.
The converse statement -- that saturation $\lambda_K^{\mathrm{mod}} =
2\pi/\beta$ characterises type-II$_\infty$ chaos -- is shown to be
\emph{false in general}: free-QFT and integrable counter-examples
(Avdoshkin--Dymarsky~\cite{AvdoshkinDymarsky2019},
Camargo--Jahnke--Kim--Nishida~\cite{Camargo2023},
Sreeram--Kannan--Modak--Aravinda~\cite{Sreeram2025})
attain the same kinematic ceiling without underlying chaos.
We therefore retain saturation only as a conjectural \emph{sufficient
condition} (Conjecture~\ref{conj:msc-saturation}).
A concrete one-sided falsifier is then derived: for the
Kolobov--Golubkov--Mu\~{n}oz~de~Nova--Steinhauer~\cite{Kolobov2021}
analog black-hole BEC at $T_H = 0.35 \pm 0.10$~nK
(Mu\~{n}oz~de~Nova et al.~\cite{Steinhauer2019}), the bound predicts
$\Gamma_{\mathrm{meas}} \le 2\pi k_B T_H/\hbar = (288 \pm 82)\,\mathrm{s}^{-1}$
for any exterior-phonon temporal correlation rate.
We compare with the concurrent boundary-area-operator construction of
Vardian~\cite{Vardian2026} (AdS/CFT) and note that our scope (de~Sitter
static patch / generic type-II$_\infty$ truncation) is complementary.
\end{abstract}

% =====================================================================
\section{Introduction}
\label{sec:intro}
% =====================================================================

The Maldacena--Shenker--Stanford (MSS) bound on chaos~\cite{MSS2016}
\begin{equation}
  \lambda_L \;\le\; \frac{2\pi k_B T}{\hbar} \;=\; \frac{2\pi}{\beta}
  \label{eq:mss}
\end{equation}
limits the OTOC Lyapunov exponent of any thermal large-$N$ quantum system,
with maximally chaotic systems (SYK, near-extremal black holes)
saturating it.
Parker, Cao, Avdoshkin, Scaffidi, and Altman~\cite{Parker2019} subsequently
introduced the universal operator-growth hypothesis: in chaotic systems
the Lanczos coefficients $b_n$ of any few-body operator grow linearly,
$b_n \sim \alpha n$, and the Krylov (Lanczos) complexity satisfies
the rigorous inequality
\begin{equation}
  \lambda_K \;\le\; 2\alpha,\qquad
  \mathcal{C}_K(t) \;\le\; A \cosh(2\alpha t),
  \label{eq:parker}
\end{equation}
sharper than~\eqref{eq:mss} and free of large-$N$ assumptions.
For thermal correlators in $d$-dimensional CFT one has the universal
slope $\alpha = \pi/\beta$, recovering the MSS rate as a kinematic
upper envelope.

In parallel, Witten~\cite{Witten2022} and Chandrasekaran--Longo--Penington--Witten~\cite{CLPW2023}
showed that gravitational observable algebras with quantum reference
frames (QRFs) are type-II$_\infty$ crossed products
$\mathcal{A}_{\mathrm{II}_\infty} = \mathcal{A}_{\mathrm{III}_1}
\rtimes_{\sigma^\phi} \mathbb{R}$, restoring a semifinite trace and
finite generalised entropy.
De~Vuyst, Eccles, H\"{o}hn, and Kirklin~\cite{DEHK2025a,DEHK2025b}
established observer-dependent gravitational entropy in this framework,
and Faulkner--Speranza~\cite{FaulknerSperanza2024} derived the
generalised second law in the crossed-product setting.

\paragraph{Aim of this paper.}
We prove a clean kinematic statement:
on \emph{any} finite-rank truncation of a type-II$_\infty$ crossed
product at the canonical Bisognano--Wichmann temperature $\beta = 2\pi$,
the modular Krylov complexity is bounded by the MSS rate
(Theorem~\ref{thm:bound}, \S\ref{sec:bound}).
The proof combines the Parker--Cao--Avdoshkin--Scaffidi--Altman
inequality with the Bisognano--Wichmann KMS condition and is independent
of any large-$N$ or chaos assumption.

\paragraph{What this paper does \emph{not} prove.}
We do \emph{not} prove the converse: saturation of the bound does
not characterise type-II$_\infty$ chaos.
Free-QFT lattice models~\cite{AvdoshkinDymarsky2019},
free and interacting bounded-power-spectrum scalar
fields~\cite{Camargo2023}, and operator/initial-state-dependent
counter-examples~\cite{Sreeram2025} all attain the same kinematic
ceiling without exhibiting MSS-saturating chaos.
The earlier ECI ``modular shadow'' draft (v1, 2026-05-04) presented
saturation as the central object; this revision (v2) downgrades
saturation to a conjecture (Conjecture~\ref{conj:msc-saturation},
\S\ref{sec:saturation}) that requires additional dynamical input
beyond the type-II$_\infty$ structure.

\paragraph{Concurrent work.}
While preparing this manuscript we became aware of
Vardian~\cite{Vardian2026} (arXiv:2602.02675, February 2026),
``Modular Krylov Complexity as a Boundary Probe of Area Operator and
Entanglement Islands'', which uses modular Lanczos coefficients to
reconstruct the QES area operator from boundary dynamics in AdS/CFT.
Vardian's analysis is complementary to ours: it operates in the
holographic AdS/CFT setting and targets the area-operator spectrum,
whereas we work intrinsically in the finite-rank type-II$_\infty$
algebra at any KMS temperature, and aim at a generic Lyapunov bound
applicable to de~Sitter and BEC analog systems.
We compare scopes in \S\ref{sec:discussion}.

\paragraph{Related work on arithmetic chaos and Krylov complexity.}
De~Clerck--Hartnoll--Santos~\cite{HartnollDCS2024}
studied the AdS black-hole interior via a Mixmaster Bianchi~IX
reduction, identifying the BKL bounce sequence with geodesic flow
on the $\Gamma(2)$ arithmetic modular surface (the same arithmetic
substrate as the Gauss-map correspondence).
Their work establishes the arithmetic chaos connection but does
\emph{not} employ a type-II$_\infty$ crossed-product completion
and does not derive a Krylov-complexity bound; our Theorem~\ref{thm:bound}
supplies the natural quantum-information completion of that programme.
In the cosmological direction, Speranza~\cite{Speranza2025} constructed
an intrinsic observer Hilbert space applicable beyond static de~Sitter,
and Requardt~\cite{Requardt2501} examined the role of type-II$_\infty$
von~Neumann algebras in quantum gravity.
Motaharfar--Siebersma--Singh~\cite{MSS2511} computed Krylov complexity
in canonical FRW (isotropic) quantum cosmology; their setting is
isotropic and does not capture the anisotropic BKL structure
studied here.

\paragraph{Literature status as of 2026-05-06.}
A systematic search of the hep-th arXiv for the period 2026-04-01 to
2026-05-06 (M5 deep literature survey, area 5: modular bootstrap/MSS)
found no published paper challenging Theorem~\ref{thm:bound} or its
three-pillar proof (Lanczos--cosh bound, Bisognano--Wichmann
K\"{a}ll\'{e}n--Lehmann spectral decomposition, inverse-Lanczos
saturation criterion; see Appendix~\ref{app:A61-proof}).
The closest 2026 works---Govindarajan--Sadanandan (2604.11277, holomorphic
MLDE rational-CFT classification) and Benjamin--Fitzpatrick--Li--Thaler
(2604.01275, machine-learning modular bootstrap at low $c$)---are
\emph{adjacent} in technique but address disjoint physics questions (RCFT
classification and spectral-gap constraints, respectively) with no
bearing on the finite-rank type-II$_\infty$ MSS bound.
The potential genus-$g \ge 2$ extension of Theorem~\ref{thm:bound} was
investigated in A77 (2026-05-05), which confirmed that the higher-genus
modular bootstrap literature lives in a different functor category (Riemann
surface partition functions / Virasoro blocks) from the type-II$_\infty$
operator-Hilbert Lanczos framework: none of the three A61 pillars has an
off-the-shelf higher-genus analog, so Theorem~\ref{thm:bound} is
\emph{not} directly extendable to $g \ge 2$ by existing methods.

\paragraph{Organization.}
\S\ref{sec:setup} fixes the algebraic setup (type-II$_\infty$ crossed
product, finite-rank truncation, modular Krylov complexity).
\S\ref{sec:bound} states and proves the bound theorem (three-step proof
sketch; full details in Appendix~\ref{app:A61-proof}).
\S\ref{sec:saturation} downgrades the saturation step to a conjecture
and reviews the free-QFT counter-examples.
\S\ref{sec:bec} applies the bound to the Kolobov--Steinhauer 2021
BEC analog black hole, deriving the falsifier
$\Gamma_{\mathrm{meas}} \le (288 \pm 82)\,\mathrm{s}^{-1}$.
\S\ref{sec:discussion} compares with Vardian and discusses scope.
\S\ref{sec:outlook} lists open problems.

% =====================================================================
\section{Setup: type-II$_\infty$ crossed product, finite-rank truncation,
         and modular Krylov complexity}
\label{sec:setup}
% =====================================================================

\subsection{The crossed-product algebra}

Let $\mathcal{A}_{\mathrm{III}_1}$ be the type-III$_1$ von~Neumann
algebra of observables in a relativistic QFT restricted to a Rindler
wedge $\mathcal{W}$, and let $\phi$ be the vacuum KMS state at the
Bisognano--Wichmann inverse temperature $\beta_{\mathrm{BW}} = 2\pi$.
Following Witten~\cite{Witten2022} and CLPW~\cite{CLPW2023}, the
crossed product with the modular automorphism group
\begin{equation}
  \mathcal{A}_{\mathrm{II}_\infty}
  \;=\;
  \mathcal{A}_{\mathrm{III}_1} \rtimes_{\sigma^\phi} \mathbb{R}
  \label{eq:crossed}
\end{equation}
is type-II$_\infty$, admits a faithful semifinite trace $\tau$, and
hosts a well-defined dressed modular Hamiltonian $K_\phi$ generating
the modular flow.
The Bisognano--Wichmann theorem~\cite{BisognanoWichmann1976} guarantees
$K_\phi = 2\pi K_{\mathrm{boost}}$ where $K_{\mathrm{boost}}$ is the
boost generator of the wedge.

\subsection{Finite-rank truncation}

Working with the full type-II$_\infty$ algebra requires care: $K_\phi$
is unbounded and the modular flow is generally non-trace-preserving.
For our bound theorem we work on the \emph{finite-rank truncation}
\begin{equation}
  \mathcal{M}_R \otimes B(\mathcal{H}_n)
  \;:=\;
  P_n \,\mathcal{A}_{\mathrm{II}_\infty}\, P_n
  \label{eq:truncation}
\end{equation}
where $P_n$ is the spectral projector of $K_\phi$ onto the modular
energy band $|K_\phi| \le \Lambda_n$ with $\dim(\mathrm{ran}\,P_n) = n$.
This truncation preserves the $\beta = 2\pi$ KMS condition (it is the
restriction of the GNS representation to a finite-dimensional subspace
on which $K_\phi$ acts as a bounded self-adjoint operator) and inherits
the semifinite trace from the type-II$_\infty$ parent.

\begin{remark}\label{rem:bw-truncation}
The Bisognano--Wichmann property survives the truncation in the
following sense: although $\mathcal{M}_R \otimes B(\mathcal{H}_n)$ is
not itself a wedge algebra, its modular flow restricted to operators
of energy $\le \Lambda_n$ agrees with the boost flow on those modes,
because the Bisognano--Wichmann modular operator commutes with the
spectral projection $P_n$.
This is the only structural fact about the truncation that we will
need.
\end{remark}

\subsection{Modular Krylov complexity}

For an operator $\mathcal{O} \in \mathcal{M}_R \otimes B(\mathcal{H}_n)$,
its modular evolution is
\begin{equation}
  \mathcal{O}(s) \;=\; e^{i K_\phi s}\,\mathcal{O}\,e^{-i K_\phi s},
  \qquad s \in \mathbb{R}.
  \label{eq:modular-flow}
\end{equation}
Equip the truncation with the KMS Gel'fand--Naimark--Segal inner product
$(A,B)_\phi := \tau(A^* B)$ (well-defined because $\tau$ is finite on
the truncation).
The Lanczos algorithm applied to the modular Liouvillian
$\mathcal{L} := [K_\phi,\,\cdot\,]$ produces an orthonormal Krylov
basis $\{|\mathcal{O}_n\rangle\}$ with tridiagonal coefficients
$\{b_n\}_{n\ge 1}$, and the modular Krylov complexity is
\begin{equation}
  \mathcal{C}_K^{\mathrm{mod}}(s)
  \;:=\;
  \sum_{n\ge 0} n\,|\varphi_n(s)|^2,
  \qquad
  \mathcal{O}(s) \;=\; \sum_{n\ge 0}\varphi_n(s)\,\mathcal{O}_n,
\end{equation}
following the modular construction of Caputa--Mag\'{a}n--Patramanis--Tonni~\cite{Caputa2024}
who showed $\lambda^{\mathrm{mod}}_L = 2\pi$ in dimensionless modular
time for 2D CFTs.

\subsection{Modular Lyapunov exponent}

The \emph{modular Lyapunov exponent} of $\mathcal{O}$ is
\begin{equation}
  \lambda_K^{\mathrm{mod}}(\mathcal{O})
  \;:=\;
  \limsup_{s\to\infty}\;\frac{1}{s}\,
  \log \mathcal{C}_K^{\mathrm{mod}}(s),
\end{equation}
when this limit is finite.
Our central object is the supremum of $\lambda_K^{\mathrm{mod}}(\mathcal{O})$
over operators $\mathcal{O}$ in a dense $\ast$-subalgebra of the
truncation; the bound theorem (next section) controls this supremum.

% =====================================================================
\section{The modular Lyapunov bound: a provable theorem}
\label{sec:bound}
% =====================================================================

\begin{theorem}[Modular Lyapunov bound on finite-rank type-II$_\infty$
truncations]\label{thm:bound}
Let $\mathcal{M}_R \otimes B(\mathcal{H}_n)$ be the finite-rank
truncation~\eqref{eq:truncation} of the type-II$_\infty$ crossed
product~\eqref{eq:crossed} at Bisognano--Wichmann KMS inverse
temperature $\beta = 2\pi$, equipped with the KMS--GNS inner product
$(A,B)_\phi = \tau(A^* B)$ and the modular flow $\sigma_s$
of~\eqref{eq:modular-flow}.
For every $\mathcal{O} \in \mathcal{M}_R \otimes B(\mathcal{H}_n)$
with finite modular two-point function, the Lanczos coefficients
$b_n(\mathcal{O})$ produced by the Lanczos algorithm with respect
to $(\,\cdot\,,\,\cdot\,)_\phi$ and Liouvillian
$\mathcal{L} = [K_\phi,\,\cdot\,]$ satisfy
\begin{equation}
  b_n(\mathcal{O}) \;\le\; n\cdot\frac{\pi}{\beta} + o(n),
  \qquad n\to\infty,
  \label{eq:bn-bound}
\end{equation}
and consequently the modular Krylov complexity satisfies
\begin{equation}
  \mathcal{C}_K^{\mathrm{mod}}(s)
  \;\le\;
  A\,\cosh\!\bigl(2 b_1(\mathcal{O}) s\bigr),
  \qquad
  \lambda_K^{\mathrm{mod}}(\mathcal{O})
  \;\le\; \frac{2\pi}{\beta} \;=\; 1,
  \label{eq:bound-main}
\end{equation}
where $A>0$ depends on $\|\mathcal{O}\|_\phi$ but not on $s$.
\end{theorem}

\begin{proof}[Proof sketch (full LMP-class proof in Appendix~\ref{app:A61-proof})]
The argument has three steps.

\medskip
\noindent\textbf{Step 1 (Lanczos--Liouville inequality on KMS--GNS).}
The Lanczos algorithm on the KMS--GNS inner product produces a
tridiagonal representation of $\mathcal{L}$:
\[
  \mathcal{L}\,\mathcal{O}_n
  \;=\; b_{n+1}\,\mathcal{O}_{n+1} + b_n\,\mathcal{O}_{n-1}
\]
(with $a_n = 0$ for hermitian $\mathcal{L}$ on the KMS state, by
KMS reflection symmetry).
Parker et al.~\cite[\S II.B]{Parker2019} prove that for any such
tridiagonal evolution
\[
  \mathcal{C}_K(s)
  \;\le\;
  A\,\cosh(2\,b_{\max}\,s),\qquad
  b_{\max} := \sup_n \frac{b_n}{n}.
\]
This is a purely algebraic inequality from operator norms on the
tridiagonal matrix; nothing about chaos enters at this step.

\medskip
\noindent\textbf{Step 2 (Bisognano--Wichmann moment bound).}
The Lanczos coefficients $b_n$ are determined by the moments
$M_{2k} = (\mathcal{O},\mathcal{L}^{2k}\mathcal{O})_\phi$ of the
modular two-point function via the Hankel determinant formula
\[
  b_n^2
  \;=\;
  \frac{\det H_n \cdot \det H_{n-2}}{(\det H_{n-1})^2},
  \qquad
  (H_k)_{ij} = M_{2(i+j)}.
\]
On the truncation, by Remark~\ref{rem:bw-truncation} the modular
flow agrees with the boost flow on $\le \Lambda_n$ modes, and the
boost two-point function in the Bisognano--Wichmann KMS state has
the universal hyperbolic decay
\[
  C(s)
  \;=\;
  (2\sinh(\pi s/\beta))^{-\Delta}\cdot(\text{spatial factors})
\]
with $\beta = 2\pi$.
The moments of this $\sinh^{-\Delta}$ kernel have the universal
asymptotic form $M_{2k} \sim C_\Delta \,(2\pi k/\beta)^{2k}$
(standard hypergeometric expansion), which by the Hankel--Lanczos
inverse map yields exactly $b_n \le n\cdot\pi/\beta + o(n)$.

\medskip
\noindent\textbf{Step 3 (combine).}
Substituting Step~2 into Step~1 gives~\eqref{eq:bound-main}.
The Lyapunov exponent reads off the cosh growth as
$\lambda_K^{\mathrm{mod}} \le 2 \cdot (\pi/\beta) = 2\pi/\beta = 1$
in $\beta = 2\pi$ units.
\end{proof}

\begin{remark}
Theorem~\ref{thm:bound} is purely kinematic.
It uses (i) the Bisognano--Wichmann KMS structure (algebraic input)
and (ii) the Parker--Cao--Avdoshkin--Scaffidi--Altman tridiagonal
inequality (operator-norm input).
No chaos hypothesis, no holography, no large-$N$ assumption is invoked.
\end{remark}

\begin{remark}\label{rem:caputa-bridge}
For 2D CFT in the Bisognano--Wichmann KMS state, Caputa--Mag\'{a}n--Patramanis--Tonni~\cite{Caputa2024}
compute $b_n^{\mathrm{CFT}} = n\cdot\pi/\beta$ exactly,
saturating~\eqref{eq:bn-bound} as an equality.
Thus the bound is tight on the CFT class.
The full type-II$_\infty$ saturation question (do all dressed
type-II$_\infty$ algebras saturate?) is open --
see~\S\ref{sec:saturation}.
\end{remark}

\begin{corollary}[MSS bound recovery]
For the truncation $\mathcal{M}_R \otimes B(\mathcal{H}_n)$ at any
KMS inverse temperature $\beta$, the modular Lyapunov exponent
satisfies $\lambda_K^{\mathrm{mod}} \le 2\pi/\beta$,
which is the MSS rate~\eqref{eq:mss}.
\end{corollary}

\begin{proof}
Replace $\beta = 2\pi$ by general $\beta$ in the proof of
Theorem~\ref{thm:bound}; the hyperbolic moment kernel becomes
$(2\sinh(\pi s/\beta))^{-\Delta}$ and the universal Lanczos slope
becomes $\pi/\beta$, giving the MSS rate.
\end{proof}

% =====================================================================
\section{Saturation as a conjecture: free-QFT counter-examples to the
         converse}
\label{sec:saturation}
% =====================================================================

The original ECI ``modular shadow'' draft~\cite{ECI_v6} proposed the
stronger statement that saturation $\lambda_K^{\mathrm{mod}} = 2\pi/\beta$
is a structural fingerprint of type-II$_\infty$.
On audit (A11, 2026-05-05), this converse fails: there exist
non-chaotic systems that saturate the kinematic ceiling.
We collect three counter-examples and then state what survives as a
conjecture.

\paragraph{Counter-example 1 (free-scalar lattice; Avdoshkin--Dymarsky 2019).}
Avdoshkin and Dymarsky~\cite{AvdoshkinDymarsky2019} show that
operator growth under imaginary-time evolution in 1D lattice models
can attain exponential Lanczos-coefficient growth even in integrable
regimes, with a 1D superexponential frequency suppression ensuring
that ``one-dimensional systems are special''.
This means linear $b_n$ and exponential $\mathcal{C}_K$ are
\emph{not} signatures of chaos; they appear in integrable lattice
systems as well.

\paragraph{Counter-example 2 (free \& interacting massive scalar QFT;
Camargo et al.\ 2023).}
Camargo, Jahnke, Kim, and Nishida~\cite{Camargo2023} compute
Krylov complexity in free and weakly interacting massive scalar
field theories and find non-chaotic regimes where the Lanczos
coefficients exhibit the same linear envelope as in maximally
chaotic systems, with mass and UV cutoff modifying the prefactor
but preserving the kinematic form.
Their work refutes the reading that linear $b_n \to 2\pi/\beta$
implies type-II$_\infty$ MSS-saturating chaos.

\paragraph{Counter-example 3 (initial-operator dependence;
Sreeram et al.\ 2025).}
Sreeram, Kannan, Modak, and Aravinda~\cite{Sreeram2025} (PRE~112
L032203) prove that Krylov complexity depends monotonically on the
inverse participation ratio (IPR) of the initial operator in the
Hamiltonian eigenbasis, and that this dependence persists in fully
chaotic regimes.
Saturation of $\lambda_K^{\mathrm{mod}}$ is therefore not a state-
or operator-independent invariant of the algebra: it varies even
within a fixed maximally-chaotic dynamics.

\paragraph{What survives as a conjecture.}
The strong claim of v1, ``saturation $\Leftrightarrow$ type-II$_\infty$'',
is refuted by counter-examples 1--3.
What remains plausible is the weaker direction: type-II$_\infty$ at
$\beta = 2\pi$ is \emph{sufficient} for the equality
$b_n = n\pi/\beta$, even if not necessary.

\begin{conjecture}[Modular saturation for dressed type-II$_\infty$
algebras]\label{conj:msc-saturation}
\emph{[CONJECTURE; v2 downgraded.]}
Let $\mathcal{A}_{\mathrm{II}_\infty}$ be a type-II$_\infty$ crossed
product~\eqref{eq:crossed} arising from a QRF-dressing of a wedge or
static-patch algebra, with the canonical $\beta = 2\pi$ KMS state.
Then on a sequence of finite-rank truncations
$\mathcal{M}_R \otimes B(\mathcal{H}_n)$, $n \to \infty$, the
Lanczos coefficients of any few-mode operator $\mathcal{O}$ in a
dense subalgebra satisfy
\[
  b_n(\mathcal{O})
  \;\xrightarrow[n\to\infty]{}\;
  n\cdot\frac{\pi}{\beta},
\]
saturating the bound~\eqref{eq:bn-bound} and yielding
$\lambda_K^{\mathrm{mod}}(\mathcal{O}) = 2\pi/\beta$ in the
large-truncation limit.
\end{conjecture}

\begin{remark}
Conjecture~\ref{conj:msc-saturation} is consistent with the
2D-CFT computation of Caputa et al.~\cite{Caputa2024}
(Remark~\ref{rem:caputa-bridge}), with the SYK Krylov saturation
of Parker et al.~\cite{Parker2019}, and with the DSSYK / sine-dilaton
holographic checks of Heller, Papalini, and
Schuhmann~\cite{HellerPapalini2024}.
We do not propose it as a theorem; the counter-examples above
show that the converse fails, and even a one-sided proof of the
sufficient direction is beyond current methods.
\end{remark}

\paragraph{Erratum: Faulkner--Speranza identification.}
The v1 draft of this work attributed to Faulkner and
Speranza~\cite{FaulknerSperanza2024} the identification
$S_{\mathrm{gen}} \leftrightarrow \mathcal{C}_K^{\mathrm{mod}}$.
On rereading, FS24 derive the crossed-product generalised entropy
and prove the GSL via Wall monotonicity; they do \emph{not}
explicitly identify $S_{\mathrm{gen}}$ with the modular Krylov
complexity.
The bridge $S_{\mathrm{gen}} \leftrightarrow \mathcal{C}_K^{\mathrm{mod}}$
is a working conjecture of the ECI programme, not a theorem of FS24.
We thank the A11/A19 audit for catching this over-attribution.

% =====================================================================
\section{Application: a falsifier on Kolobov 2021 BEC analog Hawking
         radiation}
\label{sec:bec}
% =====================================================================

We now derive a concrete observable consequence of
Theorem~\ref{thm:bound} for the analog black-hole BEC platform of
Kolobov, Golubkov, Mu\~{n}oz~de~Nova, and Steinhauer~\cite{Kolobov2021}
(Nat.~Phys.~17, 362, 2021).

\subsection{Hawking temperature and the predicted bound}

Mu\~{n}oz~de~Nova et al.~\cite{Steinhauer2019} (Nature 569, 688, 2019)
measured the analog Hawking temperature in the BEC sonic-horizon
platform as
\begin{equation}
  T_H \;=\; 0.35 \pm 0.10\,\mathrm{nK},
\end{equation}
implying the Bisognano--Wichmann KMS rate
\begin{equation}
  \boxed{\;
  \Gamma_{\mathrm{bound}}
  \;:=\;
  \frac{2\pi k_B T_H}{\hbar}
  \;=\;
  (288 \pm 82)\,\mathrm{s}^{-1}
  \;}
  \label{eq:gamma-bound}
\end{equation}
($k_B/\hbar = 1.31 \times 10^{11}\,\mathrm{rad}\,\mathrm{s}^{-1}\,\mathrm{K}^{-1}$).

By Theorem~\ref{thm:bound} applied to the (assumed) finite-rank
type-II$_\infty$ truncation of the exterior phonon algebra at the
analog horizon, any temporal correlator of exterior phonons must
have its envelope decaying with rate $\le \Gamma_{\mathrm{bound}}$:
\begin{equation}
  \Gamma_{\mathrm{meas}}
  \;\le\;
  \Gamma_{\mathrm{bound}}
  \;=\;
  (288 \pm 82)\,\mathrm{s}^{-1},
  \label{eq:gamma-pred}
\end{equation}
where $\Gamma_{\mathrm{meas}}$ is the exponential-decay rate of any
of: (i) the second-order coherence $g^{(2)}(\tau)$; (ii) the
two-point density correlator $G^{(2)}(z_1,z_2,t)$ envelope;
(iii) the Hawking-flux ramp-up rate.

\subsection{Why this is a clean falsifier}

The Kolobov~2021 dataset~\cite{Kolobov2021} contains time-resolved
density correlations at six temporal points across a 124-day,
97{,}000-repetition acquisition: the temporal envelope of
$G^{(2)}(z_1,z_2,t)$ can be extracted directly from the existing
data without new experimental runs.

\begin{proposition}[BEC bound falsifier]
If a re-analysis of Kolobov~2021 data yields
$\Gamma_{\mathrm{meas}} > 1.5 \times \Gamma_{\mathrm{bound}}
= 432\,\mathrm{s}^{-1}$ at $5\sigma$, then the type-II$_\infty$
finite-rank scaffolding underlying Theorem~\ref{thm:bound} is
falsified for the analog-BEC exterior phonon algebra.
\end{proposition}

\paragraph{Robustness to integrability counter-examples.}
Crucially, the bound~\eqref{eq:gamma-pred} is the BOUND form,
not the saturation form: it is therefore robust to the free-QFT /
integrability counter-examples
of~\S\ref{sec:saturation}~\cite{AvdoshkinDymarsky2019,Camargo2023,Sreeram2025}.
Even if the BEC phonon system is effectively integrable on the
relevant timescales, the bound still holds; only an observed
\emph{violation} would refute the type-II$_\infty$ kinematic
prefactor.

\paragraph{Cross-platform redundancy.}
Independent platforms can cross-check~\eqref{eq:gamma-pred}:
polariton-condensate analog black-hole mergers
(Solnyshkov--Paquelier--Balmisse--Malpuech 2026,
arXiv:2603.01664~\cite{Solnyshkov2026}) and water-tank vortex QNM
decay rates~\cite{TorresPatrick2017} provide independent substrates.
A UV-finite volume-law entanglement negativity
(Chandran--Fischer 2026, arXiv:2604.02075~\cite{ChandranFischer2026})
likewise yields an entanglement-growth rate bounded
by~\eqref{eq:gamma-pred} via Theorem~\ref{thm:bound}.

% =====================================================================
\section{Discussion: comparison with Vardian (2026) and scope}
\label{sec:discussion}
% =====================================================================

\paragraph{Vardian's concurrent boundary-area-operator construction.}
Vardian~\cite{Vardian2026} (arXiv:2602.02675, February 2026,
single-author hep-th) reconstructs the QES area operator from
boundary modular Lanczos coefficients in AdS/CFT, using
operator-algebra quantum error correction.
Specifically: from boundary modular dynamics,
the spectrum of the boundary modular Hamiltonian is extracted via
the Lanczos algorithm; its central component is identified with the
area operator of the dual quantum extremal surface.
Applications discussed include diagnosing island formation
and the Page transition for evaporating black holes.

\paragraph{Comparison: complementary scopes.}
Vardian's work and ours overlap in their use of modular Lanczos
coefficients on a type-II$_\infty$-like algebra, but the scopes
differ along three axes:
\begin{itemize}
\item \textbf{Holographic context.}
Vardian works in AdS/CFT and uses holographic boundary--bulk duality
explicitly; we work intrinsically in the finite-rank truncation
without invoking AdS/CFT and our results apply to any wedge or
static-patch algebra (de~Sitter, Rindler, BEC analog horizon).
\item \textbf{Object of study.}
Vardian targets the area operator (a specific element of the
modular Hamiltonian's spectrum); we target the kinematic Lyapunov
ceiling on Krylov growth (an envelope statement on all operators).
\item \textbf{Claim type.}
Vardian provides a constructive boundary--bulk dictionary;
we provide a one-sided inequality with concrete falsifier on a
laboratory analog system.
\end{itemize}
The two works together suggest that modular Krylov techniques are
becoming a standard tool for type-II$_\infty$ algebras, with
Vardian probing the holographic side and our bound theorem
covering the non-holographic side (de~Sitter, BEC, and generic
truncations).

\paragraph{Scope of Theorem~\ref{thm:bound}.}
Our theorem applies to any finite-rank truncation of a type-II$_\infty$
crossed product at canonical KMS temperature.
This covers:
\begin{enumerate}
\item Rindler wedge of a free QFT (where the bound is tight by
Caputa et al.~\cite{Caputa2024});
\item BTZ exterior in $\mathrm{AdS}_3$ (the holographic paradigm);
\item de~Sitter static patch~\cite{CLPW2023}, with $\beta = 2\pi/H$;
\item BEC and other analog-Hawking systems with a sonic horizon
(Kolobov 2021, Solnyshkov 2026);
\item the SYK / DSSYK / sine-dilaton family at large
$N$~\cite{Parker2019,HellerPapalini2024}.
\end{enumerate}
It does \emph{not} apply directly to:
\begin{enumerate}
\item Type-III$_1$ algebras at finite $N$ without a QRF dressing
(no semifinite trace, modular flow ergodic);
\item Cosmological algebras (Bianchi~IX, FRW) where the type
classification is open;
\item Type-I or II$_1$ models without boost geometry (no
Bisognano--Wichmann input).
\end{enumerate}

\paragraph{Modular vs physical Lyapunov exponents.}
The bound~\eqref{eq:bound-main} is on the modular Lyapunov exponent
$\lambda_K^{\mathrm{mod}}$, defined via modular flow.
Translating to the physical Lyapunov exponent
$\lambda_L^{\mathrm{OTOC}}$ requires the physical Hamiltonian to
align with the modular Hamiltonian under the
Connes--Rovelli~\cite{ConnesRovelli1994} thermal-time hypothesis.
For boost-symmetric wedges this alignment is automatic
(Bisognano--Wichmann); for general dynamical settings it is an
additional input.

% =====================================================================
\section{Outlook}
\label{sec:outlook}
% =====================================================================

\begin{enumerate}
\item \textbf{Lifting the truncation.}
Theorem~\ref{thm:bound} is stated on the finite-rank truncation
$\mathcal{M}_R \otimes B(\mathcal{H}_n)$.
A genuine type-II$_\infty$ statement requires a uniform-in-$n$
control of $b_n$ and a passage to the inductive-limit algebra.
The Hankel-determinant moment estimate of Step~2 in our proof
sketch is uniform in $n$ on the BW class, but the operator
$\mathcal{O}$ must lie in a dense $\ast$-subalgebra of the
type-II$_\infty$ parent.
We expect the lifted statement to hold; a complete proof is
left for future work.

\item \textbf{de~Sitter static patch.}
CLPW~\cite{CLPW2023} establish the type-II$_\infty$ structure
of the dS static patch with $\beta = 2\pi/H$ (Gibbons--Hawking
temperature).
By Theorem~\ref{thm:bound}, the modular Lyapunov exponent of
any few-mode operator on the dS static patch is bounded by
$H$ (the Hubble rate).
A concrete observable (e.g.\ a primordial-tensor-mode decoherence
rate) testing this bound would be of cosmological interest.

\item \textbf{Bianchi~IX and homogeneous chaotic cosmologies.}
The mixmaster Bianchi~IX universe exhibits chaotic dynamics with
a known modular structure on the kasner-segment wedge algebra.
Whether this algebra is type-II$_\infty$ and whether its modular
Krylov exponent saturates~\eqref{eq:bound-main} is open.

\item \textbf{Connes--St\o{}rmer entropy.}
The Connes--St\o{}rmer entropy on type-II$_1$ reductions of our
truncation provides an alternative diagnostic
(see Connes--Rovelli~\cite{ConnesRovelli1994}); whether it
inherits the same kinematic ceiling is an open question.

\item \textbf{Bridge $S_{\mathrm{gen}} \leftrightarrow
\mathcal{C}_K^{\mathrm{mod}}$.}
The bridge between the FS24 generalised entropy and modular
Krylov complexity remains a working conjecture of the ECI
programme.
A rigorous derivation would close the loop between
Theorem~\ref{thm:bound} and the Faulkner--Speranza GSL.

\item \textbf{Relation to Vardian's area-operator reconstruction.}
The boundary area operator of Vardian~\cite{Vardian2026} sits in
the central part of the modular Hamiltonian.
Its Krylov-Lanczos signature should sit at a specific level of
our truncation; making this precise would unify the two approaches.
\end{enumerate}

% =====================================================================
\section*{Acknowledgments}
% =====================================================================

This work was extracted and reworked from the ECI framework
(v6.0.53, Zenodo record \texttt{10.5281/zenodo.20030684}).
The split between the (provable) bound and the (conjectural)
saturation, and the identification of the Avdoshkin--Dymarsky /
Camargo et al.\ / Sreeram et al.\ counter-examples, owe to the A11
audit (2026-05-05).
The discovery of Vardian~\cite{Vardian2026} as concurrent
independent work was made via the same audit, with arXiv ID
2602.02675 live-verified.
Erratum on the Faulkner--Speranza over-attribution (\S\ref{sec:saturation})
follows the A11 cross-check.
We thank [collaborators] for discussions.

% =====================================================================
% APPENDIX A: Full proof of the finite-rank theorem (from A61)
% =====================================================================

\appendix

\section{Full proof of the modular Lyapunov bound on finite-rank
  type-II$_\infty$ truncations}
\label{app:A61-proof}

%% Integrated from A61_MODULAR_SHADOW_M4_FULL_PROOF/proof_draft.tex
%% (679 lines, refs live-verified 2026-05-05; M10 integration 2026-05-06).
%% All equation/theorem/lemma labels prefixed with A61: to avoid clashes.

This appendix gives the full LMP-class proof of
Theorem~\ref{thm:bound}, expanding the three-step sketch of
\S\ref{sec:bound} into a publication-quality argument.
The two main pillars are:
\begin{enumerate}
\item The universal tridiagonal-Lanczos cosh inequality of Parker,
  Cao, Avdoshkin, Scaffidi, and Altman~\cite{Parker2019} on the
  KMS--GNS inner product (algebraic, no chaos hypothesis).
\item The Bisognano--Wichmann K\"{a}ll\'{e}n--Lehmann spectral
  decomposition of the modular two-point function on the truncation
  (axiomatic QFT, no holography).
\end{enumerate}
The saturation criterion (Step~3) is a corollary of the
moment-problem inversion.

% ---- A.1 Setup and definitions ----
\subsection{Setup and definitions}
\label{ssec:A61-setup}

Throughout this appendix, $\mathcal{H}$ is a separable Hilbert space,
$\mathcal{B}(\mathcal{H})$ the bounded operators,
$\mathcal{A}_{\mathrm{III}_1}$ the type-III$_1$ wedge algebra of a
relativistic QFT in a Rindler wedge $\mathcal{W}$, and
$|\Omega\rangle$ the cyclic and separating vacuum vector inducing the
KMS state $\phi$ at Bisognano--Wichmann inverse temperature
$\beta_{\mathrm{BW}} = 2\pi$.

\subsubsection{Crossed-product algebra and modular Hamiltonian}
\label{sssec:A61-crossed}

Following Witten~\cite{Witten2022} and
Chandrasekaran--Longo--Penington--Witten (CLPW)~\cite{CLPW2023}, the
crossed product
\begin{equation}
  \mathcal{A}_{\mathrm{II}_\infty}
  \;:=\; \mathcal{A}_{\mathrm{III}_1} \rtimes_{\sigma^\phi} \mathbb{R}
  \label{eq:A61-crossed}
\end{equation}
is type-II$_\infty$, admits a faithful semifinite normal trace $\tau$,
and contains a self-adjoint generator
$K_\phi = -\beta^{-1}\log\Delta_\phi$, where $\Delta_\phi$ is the
Tomita--Takesaki modular operator of $\phi$.
The Bisognano--Wichmann theorem~\cite{BisognanoWichmann1976} identifies
$K_\phi$ on $\mathcal{A}_{\mathrm{III}_1}$ with $2\pi K_{\mathrm{boost}}$,
the wedge boost generator restricted to wedge-localised observables.
In the crossed product, $K_\phi$ extends to an unbounded self-adjoint
operator on the larger Hilbert space and generates the modular flow
\begin{equation}
  \sigma_s : \mathcal{A}_{\mathrm{II}_\infty} \to \mathcal{A}_{\mathrm{II}_\infty},
  \qquad
  \sigma_s(\mathcal{O}) \;=\; e^{i K_\phi s}\,\mathcal{O}\,e^{-i K_\phi s},
  \qquad s \in \mathbb{R}.
  \label{eq:A61-modflow}
\end{equation}

\subsubsection{Finite-rank truncation}
\label{sssec:A61-trunc}

For each $\Lambda > 0$, let $P_\Lambda$ be the spectral projector of
$K_\phi$ onto the modular-energy band $|K_\phi| \le \Lambda$, and
choose $\Lambda_n$ such that $\dim(\mathrm{ran}\,P_{\Lambda_n}) = n < \infty$.
Define the finite-rank truncation
\begin{equation}
  \mathcal{M}_R \otimes \mathcal{B}(\mathcal{H}_n)
  \;:=\;
  P_{\Lambda_n}\, \mathcal{A}_{\mathrm{II}_\infty}\, P_{\Lambda_n}.
  \label{eq:A61-trunc}
\end{equation}
This is a finite-dimensional von~Neumann algebra (an
$M_n(\mathbb{C})$-summand inside the type-II$_\infty$ parent). It
inherits the trace $\tau$ (now finite) and a faithful KMS state
$\phi_n := \tau / \tau(P_{\Lambda_n})$. Crucially, $K_\phi$ commutes
with $P_{\Lambda_n}$, so the truncated modular Hamiltonian
$K_R := P_{\Lambda_n} K_\phi P_{\Lambda_n}$ is a bounded self-adjoint
operator generating modular flow on the truncation,
$\sigma_s^{(n)}(\mathcal{O}) = e^{i K_R s}\,\mathcal{O}\,e^{-i K_R s}$.

\begin{lemma}[Bisognano--Wichmann persistence on the truncation]
\label{lem:A61-bw-trunc}
The KMS condition at $\beta = 2\pi$ holds on
$\mathcal{M}_R \otimes \mathcal{B}(\mathcal{H}_n)$ with state $\phi_n$
and modular generator $K_R$. Moreover, for any operator
$\mathcal{O} \in \mathcal{M}_R \otimes \mathcal{B}(\mathcal{H}_n)$
of modular energy support $\subset [-\Lambda_n,\Lambda_n]$, the
modular two-point function
\[
  C_\mathcal{O}(s) \;:=\; \phi_n\!\bigl(\mathcal{O}^*\,\sigma_s^{(n)}(\mathcal{O})\bigr)
\]
agrees on $|s| \le s_n$ (with $s_n \to \infty$ as $\Lambda_n\to\infty$)
with the boost two-point function on the wedge, i.e.\ with the analytic
continuation of the Wightman two-point function at boost angle $s$.
\end{lemma}

\begin{proof}[Sketch]
The KMS condition is preserved under spectral truncation because
$\Delta_\phi P_{\Lambda_n} = P_{\Lambda_n}\Delta_\phi$ implies
$\sigma_s$ restricts to $\sigma_s^{(n)}$ on $\mathrm{ran}(P_{\Lambda_n})$.
The Bisognano--Wichmann identification $K_\phi = 2\pi K_{\mathrm{boost}}$
descends to $K_R = 2\pi P_{\Lambda_n} K_{\mathrm{boost}} P_{\Lambda_n}$;
the boost two-point function inherits this structure on the
modular-energy band $\le \Lambda_n$. The cutoff horizon
$s_n = O(\log \Lambda_n)$ controls when spectral truncation begins to
deform the analytic two-point function.
\end{proof}

\subsubsection{Modular Krylov complexity}
\label{sssec:A61-krylov}

Equip $\mathcal{M}_R \otimes \mathcal{B}(\mathcal{H}_n)$ with the
KMS--GNS inner product
\begin{equation}
  (A,B)_{\phi_n} \;:=\; \tau\!\bigl(\rho^{1/2} A^* \rho^{1/2} B\bigr)
  \;=\; \langle \Omega_n | A^* \sigma_{i\beta/2}^{(n)}(B) |\Omega_n \rangle,
  \label{eq:A61-gns}
\end{equation}
where $\rho := e^{-\beta K_R}/Z$ is the truncated KMS density operator
and $|\Omega_n\rangle$ the cyclic vector. The modular Liouvillian
\begin{equation}
  \mathcal{L} := i [K_R,\,\cdot\,]
  \label{eq:A61-liouv}
\end{equation}
acts on the operator Hilbert space
$(\mathcal{M}_R \otimes \mathcal{B}(\mathcal{H}_n), (\cdot,\cdot)_{\phi_n})$
as a self-adjoint generator of $\sigma_s^{(n)}$. (Self-adjointness
follows from KMS reflection symmetry
$(A, [K_R, B])_{\phi_n} = ([K_R, A], B)_{\phi_n}^*$; we use the
Hermiticity convention with $i$ explicit so that
$\mathcal{L} = \mathcal{L}^\dagger$.)

Given $\mathcal{O} \in \mathcal{M}_R \otimes \mathcal{B}(\mathcal{H}_n)$
with $(\mathcal{O},\mathcal{O})_{\phi_n} = 1$, the Lanczos algorithm
produces an orthonormal basis $\{|\mathcal{O}_n\rangle\}_{n\ge 0}$ of
the Krylov subspace
$\overline{\mathrm{span}}\{|\mathcal{O}\rangle, \mathcal{L}|\mathcal{O}\rangle, \mathcal{L}^2|\mathcal{O}\rangle, \ldots\}$
satisfying the three-term recurrence
\begin{equation}
  \mathcal{L}\,|\mathcal{O}_n\rangle
  \;=\; b_{n+1}\,|\mathcal{O}_{n+1}\rangle + b_n\,|\mathcal{O}_{n-1}\rangle,
  \qquad
  b_n > 0,\quad |\mathcal{O}_{-1}\rangle := 0,
  \label{eq:A61-lanczos}
\end{equation}
with $a_n \equiv 0$ by KMS reflection symmetry of $\mathcal{L}$ on
$(\cdot,\cdot)_{\phi_n}$. The modular Krylov complexity is
\begin{equation}
  \mathcal{C}_K^{\mathrm{mod}}(s)
  \;:=\;
  \sum_{n \ge 0} n\,|\varphi_n(s)|^2,
  \qquad
  e^{i\mathcal{L}s}|\mathcal{O}\rangle
  \;=\;
  \sum_{n \ge 0} \varphi_n(s)\,|\mathcal{O}_n\rangle.
  \label{eq:A61-CKmod}
\end{equation}

% ---- A.2 Full three-step proof ----
\subsection{Full proof of Theorem~\ref{thm:bound}}
\label{ssec:A61-proof}

The proof has three steps: (i) the Parker--Cao--Avdoshkin--Scaffidi--Altman
tridiagonal-Lanczos inequality on the KMS--GNS inner product (operator-algebraic,
no chaos); (ii) the Bisognano--Wichmann moment estimate via the Hankel
determinant (kinematic, no holography); (iii) the inverse-Lanczos
saturation criterion (reduction of equality to the two-point function form,
with free-QFT counter-examples precluding a chaos reading).

\medskip

\noindent\textbf{Step 1: Tridiagonal-Lanczos cosh inequality on the
KMS--GNS inner product.}

Let $\mathcal{O} \in \mathcal{M}_R \otimes \mathcal{B}(\mathcal{H}_n)$
with $(\mathcal{O},\mathcal{O})_{\phi_n} = 1$. Let
$\mathcal{O}(s) = e^{i\mathcal{L}s}\mathcal{O}$ and decompose
$\mathcal{O}(s) = \sum_n \varphi_n(s)\,|\mathcal{O}_n\rangle$
in the Krylov basis~\eqref{eq:A61-lanczos}.
The amplitudes $\varphi_n(s)$ satisfy the discrete Schr\"{o}dinger equation
\begin{equation}
  i\,\partial_s \varphi_n(s)
  \;=\;
  b_{n+1}\,\varphi_{n+1}(s) + b_n\,\varphi_{n-1}(s),
  \qquad
  \varphi_n(0) = \delta_{n,0}.
  \label{eq:A61-schrod}
\end{equation}

\emph{Sub-step 1a (energy conservation).} The norm
$\sum_n |\varphi_n(s)|^2 = 1$ is conserved (unitarity of $e^{i\mathcal{L}s}$
on the KMS--GNS Hilbert space).

\emph{Sub-step 1b (operator norm of the truncated tridiagonal generator).}
For the truncated tridiagonal matrix $T^{(N)}$ with entries
$T^{(N)}_{n,n+1} = T^{(N)}_{n+1,n} = b_{n+1}$ and zero elsewhere,
the standard Gershgorin bound gives
$\|T^{(N)}\|_{\mathrm{op}} \le 2 \max_{1\le n\le N} b_n$.

\emph{Sub-step 1c (Parker et al.~\cite{Parker2019} cosh envelope).}
Following Parker--Cao--Avdoshkin--Scaffidi--Altman~\cite{Parker2019}
on a tridiagonal Hamiltonian with linearly bounded coefficients
$b_n \le \alpha n$, the modular Krylov complexity satisfies
\begin{equation}
  \mathcal{C}_K^{\mathrm{mod}}(s)
  \;\le\;
  A\,\cosh(2\alpha s)
  \label{eq:A61-parker-cosh}
\end{equation}
for some constant $A$ depending on $b_1$ but not on $s$.
The proof in~\cite{Parker2019} maps the recurrence~\eqref{eq:A61-schrod}
to a chiral Dirac equation on a half-line and applies a
propagation-speed (Lieb--Robinson) bound, yielding the cosh growth as
the saturating profile.

Adapting to the modular setting: replace the physical Hamiltonian $H$
by $K_R$, the physical Liouvillian $[H,\,\cdot\,]$ by the modular
Liouvillian $\mathcal{L} = i[K_R,\,\cdot\,]$, and the Gibbs inner
product $\mathrm{tr}(e^{-\beta H} A^* B)$ by the truncated KMS--GNS inner
product~\eqref{eq:A61-gns}. The proof of~\eqref{eq:A61-parker-cosh}
in~\cite{Parker2019} uses only the tridiagonal structure of the
recurrence, the Hermiticity of the generator, and the unitarity of
evolution; none of these properties are altered by substituting
$\mathcal{L}$ for the physical Liouvillian.

Therefore: \emph{if} $b_n \le \alpha n$, then~\eqref{eq:A61-parker-cosh}
gives $\lambda_K^{\mathrm{mod}}(\mathcal{O}) \le 2\alpha$.
This reduces Theorem~\ref{thm:bound} to establishing
$\alpha \le \pi/\beta$, which is Step~2.

\medskip

\noindent\textbf{Step 2: Bisognano--Wichmann moment bound and the
universal $\pi/\beta$ slope.}

The modular Lanczos coefficients $b_n^2$ are determined by the moments
\begin{equation}
  M_{2k} := (\mathcal{O}, \mathcal{L}^{2k}\mathcal{O})_{\phi_n}
  \;=\; (-1)^k C_\mathcal{O}^{(2k)}(0),
  \label{eq:A61-moments}
\end{equation}
i.e.\ by the even derivatives of the modular two-point function
$C_\mathcal{O}(s) := \phi_n(\mathcal{O}^* \sigma_s^{(n)}(\mathcal{O}))$
at $s=0$. The Hankel--Lanczos relation reads
\begin{equation}
  b_n^2 \;=\;
  \frac{D_n \cdot D_{n-2}}{(D_{n-1})^2},
  \qquad
  D_k := \det H_k,\quad
  (H_k)_{ij} = M_{2(i+j)},\;\; 0 \le i,j \le k.
  \label{eq:A61-hankel}
\end{equation}
This is the standard Stieltjes--Chebyshev moment-problem formula;
see Viswanath--M\"{u}ller, \emph{The Recursion Method} (1994).

\emph{Sub-step 2a (Bisognano--Wichmann spectral form on the truncation).}
By Lemma~\ref{lem:A61-bw-trunc}, on $|s|\le s_n$ the modular
two-point function agrees with the boost two-point function of a
wedge-localised operator. By Bisognano--Wichmann~\cite{BisognanoWichmann1976},
this function in the KMS state at $\beta = 2\pi$ admits a K\"{a}ll\'{e}n--Lehmann
spectral decomposition over scaling-dimension weights:
\begin{equation}
  C_\mathcal{O}(s)
  \;=\;
  \int_0^\infty \!d\Delta\;\rho_\mathcal{O}(\Delta)
  \,\bigl(2\sinh(\pi s/\beta)\bigr)^{-2\Delta}
  + (\text{spatial form factor}),
  \label{eq:A61-bw-spectral}
\end{equation}
where $\rho_\mathcal{O}(\Delta) \ge 0$ is the operator-dependent
weight on conformal dimensions $\Delta$. This is the analogue of the
K\"{a}ll\'{e}n--Lehmann representation under modular flow, well-known in
the algebraic-QFT literature; see Haag~\cite{Haag1996} (Chapter V)
and Borchers' modular reflection identity.
For the truncation, the integral~\eqref{eq:A61-bw-spectral} is cut off
at the maximal scaling dimension $\Delta_n = O(\Lambda_n)$ accessible
to the modular-energy band.

\emph{Sub-step 2b (moments of a hyperbolic kernel).}
For each spectral weight $(2\sinh(\pi s/\beta))^{-2\Delta}$, the
Taylor expansion at $s=0$ involves only even powers, and the moments
\[
  M_{2k}^{(\Delta)} := (-1)^k \frac{d^{2k}}{ds^{2k}}
  \bigl(2\sinh(\pi s/\beta)\bigr)^{-2\Delta}\Big|_{s=0}
\]
admit the closed-form asymptotic
\begin{equation}
  M_{2k}^{(\Delta)}
  \;\sim\;
  C_\Delta \,(2\pi k/\beta)^{2k}
  \,\bigl(1 + O(1/k)\bigr),
  \qquad k\to\infty,
  \label{eq:A61-moment-asymp}
\end{equation}
with $C_\Delta > 0$ depending on $\Delta$ but not on $k$.
This follows from the saddle-point evaluation of the inverse Mellin
transform; see Whittaker--Watson, \emph{Modern Analysis} (1927) \S 13.6,
or directly by integrating the spectral representation
$(2\sinh(\pi s/\beta))^{-2\Delta} = \int_{-\infty}^\infty e^{i\omega s}
\tilde\rho_\Delta(\omega)\,d\omega$ where $\tilde\rho_\Delta(\omega)$
is peaked near $|\omega| \sim 2\pi k/\beta$ for the $2k$-th moment.

Integrating against $\rho_\mathcal{O}(\Delta)$ gives
\begin{equation}
  M_{2k} \;\le\; C_\mathcal{O}\,(2\pi k/\beta)^{2k}
  \,\bigl(1 + O(1/k)\bigr),
  \qquad k\to\infty,
  \label{eq:A61-total-moment}
\end{equation}
where $C_\mathcal{O} := \int_0^{\Delta_n}d\Delta\,\rho_\mathcal{O}(\Delta)
\,C_\Delta < \infty$ on the truncation.

\emph{Sub-step 2c (Hankel--Lanczos inversion).}
Substituting~\eqref{eq:A61-total-moment} into the Hankel
formula~\eqref{eq:A61-hankel} and using the standard result (Akhiezer,
\emph{Classical Moment Problem}, 1965, \S 1.4) that for moments
$M_{2k} = m^{2k}\,\mu_k$ with $\mu_k$ subexponential the corresponding
Lanczos coefficients satisfy $b_n = m\cdot n + o(n)$, one obtains
\begin{equation}
  b_n(\mathcal{O}) \;\le\; n\cdot \frac{\pi}{\beta} + o(n),
  \qquad n \to \infty,
\end{equation}
which is~\eqref{eq:bn-bound}. Positivity of the Hankel determinants
is inherited from the positive-semidefinite spectral weight
$\rho_\mathcal{O} \ge 0$.

\medskip

\noindent\textbf{Step 3: Saturation criterion and the universal
operator-growth hypothesis (UOGH).}

Equality $b_n(\mathcal{O}) = n\pi/\beta$ holds asymptotically if and
only if~\eqref{eq:A61-moment-asymp} saturates, i.e.\ the spectral
weight $\rho_\mathcal{O}$ is concentrated on a single dimension or
yields equivalent moment growth. By the inverse-Lanczos statement of
Parker et al.~\cite[Theorem~2]{Parker2019}, specifying the asymptotic
slope $\alpha = \pi/\beta$ uniquely determines the long-time exponential
decay as $C_\mathcal{O}(s) \sim e^{-(\pi/\beta) s\cdot 2\Delta}$,
equivalent to
\begin{equation}
  C_\mathcal{O}(s)
  \;=\;
  N_\Delta\,\bigl(2\sinh(\pi s/\beta)\bigr)^{-\Delta},
  \label{eq:A61-bw-twopt-sat}
\end{equation}
the Bisognano--Wichmann form. This is the UOGH of Parker et
al.~\cite{Parker2019}, in its modular incarnation.

The chaos reading is refuted by free-QFT counter-examples
(see \S\ref{sec:saturation} for full discussion):
Avdoshkin--Dymarsky~\cite{AvdoshkinDymarsky2019} show linear $b_n$
in integrable lattice models; Camargo--Jahnke--Kim--Nishida~\cite{Camargo2023}
show the same kinematic form in free and weakly-interacting massive
scalar QFTs; Sreeram--Kannan--Modak--Aravinda~\cite{Sreeram2025}
show saturation is operator-dependent within maximally-chaotic dynamics.

\medskip

\noindent\textbf{Combining Steps 1, 2, 3.}

Step~2 gives $b_n(\mathcal{O}) \le n\cdot\pi/\beta + o(n)$, so
$\alpha := \pi/\beta$ is a valid linear bound. Step~1 then yields
$\mathcal{C}_K^{\mathrm{mod}}(s) \le A\cosh(2\pi s/\beta)$, giving
$\lambda_K^{\mathrm{mod}}(\mathcal{O}) \le 2\pi/\beta$.
At $\beta = 2\pi$ this is the MSS rate $\lambda_K^{\mathrm{mod}} \le 1$.
Step~3 characterises saturation~\eqref{eq:A61-bw-twopt-sat} as a
kinematic, not dynamical, statement, with free-QFT counter-examples
precluding a chaos reading. \qed

% ---- A.3 Concurrent independent work (comparison) ----
\subsection{Comparison with Vardian (2026)}
\label{ssec:A61-vardian}

Vardian~\cite{Vardian2026} (arXiv:2602.02675, February 2026,
single-author hep-th) presents a complementary construction in AdS/CFT.
The central object is the boundary modular Hamiltonian's spectrum,
extracted via the Lanczos algorithm and identified through
operator-algebra quantum error correction (OAQEC) with the area
operator of the dual quantum extremal surface (QES).

Our Theorem~\ref{thm:bound} and Vardian's construction share the
modular-Lanczos technical machinery but address disjoint physical
questions. Vardian operates in AdS/CFT and targets a specific spectral
element (the area operator); we target the kinematic Lyapunov ceiling on
Krylov growth for all operators in any finite-rank type-II$_\infty$
truncation, including non-holographic settings (de~Sitter, BEC).
The two works together confirm that modular Krylov techniques are a
productive tool for type-II$_\infty$ algebras.

% ---- A.4 Open problems ----
\subsection{Open problems from the full proof}
\label{ssec:A61-open}

\begin{enumerate}
\item \textbf{Lifting to the full type-II$_\infty$ algebra.}
Theorem~\ref{thm:bound} is on the finite-rank truncation.
Passing to the inductive-limit algebra requires uniform-in-$n$ control
of $b_n$ for operators in a dense $\ast$-subalgebra of
$\mathcal{A}_{\mathrm{II}_\infty}$.

\item \textbf{Saturation as a sufficient (not necessary) condition.}
Conjecture~\ref{conj:msc-saturation} (\S\ref{sec:saturation}) remains
open.

\item \textbf{Bridge $S_{\mathrm{gen}} \leftrightarrow
\mathcal{C}_K^{\mathrm{mod}}$.}
The Faulkner--Speranza~\cite{FaulknerSperanza2024} generalised entropy
and $\mathcal{C}_K^{\mathrm{mod}}$ are conjecturally related; rigorous
statement is open.

\item \textbf{Higher-genus extension.}
As confirmed by the A77 investigation (2026-05-05), the genus-$g \ge 2$
modular bootstrap lives in a different functor category (Riemann surface
partition functions / Virasoro blocks) from the type-II$_\infty$
operator-Hilbert Lanczos framework of this appendix: the three pillars of
Step~1--3 have no off-the-shelf higher-genus analogs. Extension to $g \ge 2$
is a genuine open problem requiring new mathematical technology.
\end{enumerate}

% =====================================================================
\begin{thebibliography}{99}
% =====================================================================

\bibitem{MSS2016}
J.~Maldacena, S.~H.~Shenker, and D.~Stanford,
``A bound on chaos,''
\textit{JHEP} \textbf{08} (2016) 106,
\href{https://arxiv.org/abs/1503.01409}{arXiv:1503.01409 [hep-th]}.

\bibitem{Parker2019}
D.~E.~Parker, X.~Cao, A.~Avdoshkin, T.~Scaffidi, and E.~Altman,
``A Universal Operator Growth Hypothesis,''
\textit{Phys.\ Rev.\ X} \textbf{9} (2019) 041017,
\href{https://arxiv.org/abs/1812.08657}{arXiv:1812.08657 [cond-mat.str-el]}.

\bibitem{Caputa2024}
P.~Caputa, J.~M.~Mag\'{a}n, D.~Patramanis, and E.~Tonni,
``Krylov complexity of modular Hamiltonian evolution,''
\textit{Phys.\ Rev.\ D} \textbf{109} (2024) 086004,
\href{https://arxiv.org/abs/2306.14732}{arXiv:2306.14732 [hep-th]}.

\bibitem{Witten2022}
E.~Witten,
``Gravity and the crossed product,''
\textit{JHEP} \textbf{10} (2022) 008,
\href{https://arxiv.org/abs/2112.12828}{arXiv:2112.12828 [hep-th]}.

\bibitem{CLPW2023}
V.~Chandrasekaran, R.~Longo, G.~Penington, and E.~Witten,
``An algebra of observables for de~Sitter space,''
\textit{JHEP} \textbf{02} (2023) 082,
\href{https://arxiv.org/abs/2206.10780}{arXiv:2206.10780 [hep-th]}.

\bibitem{FaulknerSperanza2024}
T.~Faulkner and A.~J.~Speranza,
``Gravitational algebras and the generalized second law,''
\textit{JHEP} \textbf{11} (2024) 099,
\href{https://arxiv.org/abs/2405.00847}{arXiv:2405.00847 [hep-th]}.
%% v2 erratum: FS24 derives crossed-product S_gen and the GSL via Wall monotonicity;
%% it does NOT explicitly identify S_gen with modular Krylov complexity.
%% The bridge S_gen ↔ C_K^mod is an ECI working conjecture, not a FS24 theorem.

\bibitem{DEHK2025a}
J.~De~Vuyst, S.~Eccles, P.~A.~H\"{o}hn, and J.~Kirklin,
``Crossed products and quantum reference frames: on the observer-dependence
of gravitational entropy,''
\textit{JHEP} \textbf{07} (2025) 063,
\href{https://arxiv.org/abs/2412.15502}{arXiv:2412.15502 [hep-th]}.

\bibitem{DEHK2025b}
J.~De~Vuyst, S.~Eccles, P.~A.~H\"{o}hn, and J.~Kirklin,
``Gravitational entropy is observer-dependent,''
\textit{JHEP} \textbf{07} (2025) 146,
\href{https://arxiv.org/abs/2405.00114}{arXiv:2405.00114 [hep-th]}.

\bibitem{HellerPapalini2024}
M.~P.~Heller, A.~Papalini, and T.~Schuhmann,
``Krylov spread complexity as holographic complexity beyond JT gravity,''
\textit{Phys.\ Rev.\ Lett.}\ \textbf{135} (2025) 151602,
\href{https://arxiv.org/abs/2412.17785}{arXiv:2412.17785 [hep-th]}.

\bibitem{BisognanoWichmann1976}
J.~J.~Bisognano and E.~H.~Wichmann,
``On the duality condition for a Hermitian scalar field,''
\textit{J.\ Math.\ Phys.}\ \textbf{17} (1976) 303.

\bibitem{Haag1996}
R.~Haag, \textit{Local Quantum Physics: Fields, Particles, Algebras},
2nd ed., Springer (1996).
%% Chapter V: K\"all\'en--Lehmann representation under modular flow,
%% used in Appendix~\ref{app:A61-proof} Step 2 (Sub-step 2a).

\bibitem{ConnesRovelli1994}
A.~Connes and C.~Rovelli,
``Von Neumann algebra automorphisms and time-thermodynamics relation in
generally covariant quantum theories,''
\textit{Class.\ Quantum Grav.}\ \textbf{11} (1994) 2899--2917.

%% --- Vardian (concurrent independent work, must cite) ---
\bibitem{Vardian2026}
N.~Vardian,
``Modular Krylov Complexity as a Boundary Probe of Area Operator and
Entanglement Islands,''
\href{https://arxiv.org/abs/2602.02675}{arXiv:2602.02675 [hep-th]}
(February 2026).
%% Single-author hep-th. Live-verified 2026-05-05.

%% --- Free-QFT counter-examples to the saturation converse ---
\bibitem{AvdoshkinDymarsky2019}
A.~Avdoshkin and A.~Dymarsky,
``Euclidean operator growth and quantum chaos,''
\textit{Phys.\ Rev.\ Research} \textbf{2} (2020) 043234,
\href{https://arxiv.org/abs/1911.09672}{arXiv:1911.09672 [cond-mat.stat-mech]}.

\bibitem{Camargo2023}
H.~A.~Camargo, V.~Jahnke, K.-Y.~Kim, and M.~Nishida,
``Krylov Complexity in Free and Interacting Scalar Field Theories with
Bounded Power Spectrum,''
\textit{JHEP} \textbf{05} (2023) 226,
\href{https://arxiv.org/abs/2212.14702}{arXiv:2212.14702 [hep-th]}.

\bibitem{Sreeram2025}
Sreeram~PG, J.~Bharathi~Kannan, R.~Modak, and S.~Aravinda,
``Dependence of Krylov complexity on the initial operator and state,''
\textit{Phys.\ Rev.\ E} \textbf{112} (2025) L032203,
\href{https://arxiv.org/abs/2503.03400}{arXiv:2503.03400 [quant-ph]}.

%% --- BEC analog Hawking platforms (A19 BOUND falsifier) ---
\bibitem{Kolobov2021}
V.~I.~Kolobov, K.~Golubkov, J.~R.~Mu\~{n}oz~de~Nova, and J.~Steinhauer,
``Observation of stationary spontaneous Hawking radiation and the time
evolution of an analogue black hole,''
\textit{Nat.\ Phys.}\ \textbf{17} (2021) 362--367,
\href{https://arxiv.org/abs/1910.09363}{arXiv:1910.09363 [gr-qc]}.

\bibitem{Steinhauer2019}
J.~R.~Mu\~{n}oz~de~Nova, K.~Golubkov, V.~I.~Kolobov, and J.~Steinhauer,
``Observation of thermal Hawking radiation at the Hawking temperature in
an analogue black hole,''
\textit{Nature} \textbf{569} (2019) 688--691,
\href{https://arxiv.org/abs/1809.00913}{arXiv:1809.00913 [gr-qc]}.

\bibitem{Solnyshkov2026}
D.~D.~Solnyshkov, I.~Paquelier, A.~Balmisse, and G.~Malpuech,
``Polariton-condensate analog black-hole merger'' (working title),
\href{https://arxiv.org/abs/2603.01664}{arXiv:2603.01664}
(2026).

\bibitem{ChandranFischer2026}
A.~Chandran and U.~R.~Fischer,
``UV-finite volume-law entanglement negativity'' (working title),
\href{https://arxiv.org/abs/2604.02075}{arXiv:2604.02075}
(2026).

\bibitem{TorresPatrick2017}
T.~Torres, S.~Patrick, A.~Coutant, M.~Richartz, E.~W.~Tedford, and S.~Weinfurtner,
``Rotational superradiant scattering in a vortex flow,''
\textit{Nat.\ Phys.}\ \textbf{13} (2017) 833--836.

%% --- ECI internal references ---
\bibitem{ECI_v6}
[ECI author(s)],
``Entanglement--Complexity--Information: a phenomenological architecture
for quantum gravity phenomenology (v6.0.53),''
Zenodo record \texttt{10.5281/zenodo.20030684} (2026).

%%  R1 related-work references (M59, 2026-05-06; all arXiv IDs live-verified)
\bibitem{HartnollDCS2024}
M.~De~Clerck, S.~A.~Hartnoll, and J.~E.~Santos,
``Mixmaster chaos in an AdS black hole interior,''
\textit{JHEP} \textbf{07} (2024) 202,
\href{https://arxiv.org/abs/2312.11622}{arXiv:2312.11622 [hep-th]}.
%% M48 verified JHEP07(2024)202; author order live-verified 2026-05-06.

\bibitem{Requardt2501}
M.~Requardt,
``The Role of Type~II$_\infty$ v.Neumann Algebras and their Tensor Structure
  in Quantum Gravity,''
\href{https://arxiv.org/abs/2501.06009}{arXiv:2501.06009 [gr-qc]} (2025).

\bibitem{MSS2511}
M.~Motaharfar, M.~R.~Siebersma, and P.~Singh,
``Krylov Complexity in Canonical Quantum Cosmology,''
\href{https://arxiv.org/abs/2511.17711}{arXiv:2511.17711} (2025).
%% FRW (isotropic) setting; distinct from Bianchi~IX anisotropic case.

\end{thebibliography}

\end{document}
